\section{The Winter Ember Days: Annunciation, Visitation, and Harvest of Mercy}\label{sec:ember-days}

Between the joy of Gaudete Sunday and the steep final climb toward Christmas
stand three days unlike any other weekdays of Advent: the Wednesday, Friday,
and Saturday after the Third Sunday. The Roman books call them \textit{Quatuor
Tempora Adventus}, the Winter Ember Days, and rank them as second-class ferias
with Masses and Gospel homilies of their own. Depending on when the First
Sunday of Advent opens, Ember Wednesday falls between 14 and 20 December, Ember
Friday between 16 and 22 December, and Ember Saturday between 17 and 23
December---so that in every year, Ember Saturday already stands inside the week
of the Great Antiphons.

Long before Rome called this season Advent, the popes gathered the city in the
tenth month---December, counting from March---to thank God for the completed
harvest of the year with fasting, prayer, and almsgiving. St.~Leo the Great
sets the reason before his congregation with crisp plainness (Sermon~17 on the
Fast of the Tenth Month):

\begin{quote}
Among these is also the solemn fast of the tenth month, which is now to be kept
by us according to yearly custom, because it is altogether just and godly to
give thanks to the Divine bounty for the crops which the earth has produced for
the use of men under the guiding hand of supreme Providence. And to show that
we do this with ready mind, we must exercise not only the self-restraint of
fasting, but also diligence in almsgiving, that from the ground of our heart
also may spring the germ of righteousness and the fruit of love, and that we
may deserve God's mercy by showing mercy to His poor.
\end{quote}

In Sermon 12, Leo calls his flock to the three days by name: ``On Wednesday and
Friday therefore let us fast: and on Saturday let us keep vigil with the most
blessed Apostle Peter, who will deign to aid our supplications and fast and
alms with his own prayers.'' By the time the Roman Missal reached its classic
shape, those three days of winter fasting had been woven into a quiet triptych
of the Incarnation itself: the Annunciation on Wednesday, the Visitation on
Friday, and the long ordination vigil of Saturday beside St.~Peter's tomb.

\subsection{Ember Wednesday: The Golden Mass at St.~Mary Major}

On Wednesday the Roman station returns to St.~Mary Major, and the Mass opens
with the cry of Isaiah 45:8, \textit{Rorate caeli desuper}, ``Drop down dew, ye
heavens, from above, and let the clouds rain the just: let the earth be opened,
and bud forth a saviour.'' Because the Gospel is St.~Luke's narrative of the
Annunciation (Luke 1:26--38), medieval devotion knew this liturgy as the
\textit{Missa Aurea}, the Golden Mass.

Unlike an ordinary weekday, Ember Wednesday hears two Old Testament lessons,
each preceded by a Collect and followed by a Gradual. First Isaiah 2:2--5
promises that all nations will flow to the mountain of the Lord's house; then,
before the second lesson, the Church prays \textit{Festina, quaesumus, Domine},
pleading that the Lord delay no longer:

\begin{massbox}{Second Collect of Ember Wednesday}
\textbf{Latin.} Festina, quaesumus, Domine, ne tardaveris, et auxilium nobis
supernae virtutis impende: ut adventus tui consolationibus subleventur, qui in
tua pietate confidunt: Qui vivis et regnas cum Deo Patre in unitate.

\textbf{English.} Hasten, we beseech Thee, O Lord, do not tarry, and grant us
the help of Thy heavenly power : that those who trust in Thy goodness may be
helped by the consolations of Thy coming : Who livest.
\end{massbox}

Then Isaiah 7:10--15 speaks the sign given to the house of David: ``Behold a
virgin shall conceive, and bear a son and his name shall be called Emmanuel.''
The Gospel immediately meets the prophet's voice with Gabriel entering a town
in Galilee named Nazareth. At Matins on this day, all three lessons are drawn
from St.~Ambrose's exposition of Luke 1:26, explaining why God willed the
Mother of the Saviour to be espoused to Joseph before the Spirit overshadowed
her: ``Perhaps it was lest any might blasphemously say that she had conceived
in adultery the Holy One.''

No passage in Christian preaching catches the suspense of that moment more
vividly than St.~Bernard of Clairvaux in his Fourth Homily on the
\textit{Missus est}. Everything created---patriarchs in the shades of death,
angels in heaven, the whole exiled race of Adam---hangs in silence upon one
girl's answer in a small house of Galilee:

\begin{quote}
Adam, with all his race---Adam, a weeping exile from Paradise, implores it of
you. Abraham entreats you, David beseeches you. This is the object of the
burning desires of the holy fathers, of your fathers, who are still dwelling in
the region of the shades of death. Behold the entire human race prostrate at
your feet in expectation. And rightly, for on your word depend the consolation
of the wretched, the redemption of the captive, the freedom of the condemned,
the salvation of your entire race, of all the children of Adam. Hasten, then, O
Lady, to give your answer; hasten to speak the word so longed for by all on
earth, in limbo, and in heaven. Yea, the King and Lord of all things, Who has
greatly desired your beauty, desires as eagerly your word of consent, by which
He has purposed to save the world.
\end{quote}

A few lines later Bernard presses his appeal: ``Say the word and receive the
Word. Utter thy human word and conceive the Divine Word. Pronounce the
transitory word and embrace the Word everlasting.'' St.~Thomas Aquinas states
the theology beneath Bernard's poetry with stark precision in the
\textit{Summa Theologiae} (III, q.~30, a.~1): the Annunciation took place ``in
order to show that there is a certain spiritual wedlock between the Son of God
and human nature. Wherefore in the Annunciation the Virgin's consent was
besought in lieu of that of the entire human nature.'' God would not force our
redemption upon us without asking a human yes; and Mary gave it: ``Behold the
handmaid of the Lord: be it done to me according to thy word'' (Luke 1:38).

\subsection{Ember Friday: Haste Into the Hill Country}

On Friday the station moves to the Basilica of the Twelve Holy Apostles. The
Introit, \textit{Prope es tu, Domine} (``Thou art near, O Lord,'' Ps~118~[119]:151),
speaks to One who is no longer far off. The lesson is Isaiah 11:1--5: ``And
there shall come forth a rod out of the root of Jesse, and a flower shall rise
up out of his root.'' Immediately the Gospel carries us up the road from
Nazareth into the hill country of Judea, where Mary greets Elizabeth and the
unborn John leaps in his mother's womb (Luke 1:39--47).

At Ember Friday Matins, St.~Ambrose (\textit{Exposition of the Gospel of Luke},
Book~II) watches Mary leave Nazareth with haste and refuses to let anyone read
her hurry as doubt or curiosity:

\begin{quote}
When the holy Virgin had heard it, she arose and went to visit her cousin. She
did not go to see if what she had heard was true, because she did not believe
God, or because she knew not who the messenger had been, or yet because she
doubted the fact adduced in proof. She went joyfully as one who hath received a
mercy in answer to his vow goeth to pay the same. She went with devotion, as a
godly person goeth to execute a religious duty. She went into the hill country
in joyful haste. And is it not something that she went up into the hills? God
was already in her womb, and her feeling bore her continually upward. The grace
of the Holy Spirit knoweth no slow working.
\end{quote}

In the same commentary St.~Ambrose adds a sentence that turns the Visitation
into a prayer for every Christian household: ``If according to the flesh there
is one mother of Christ, yet, according to faith, Christ is the fruit of all.
For every soul receives the word of God if only he be unspotted and free from
sin, and preserves it with unsullied purity.'' And St.~Bede notes why Mary
sings, ``My spirit hath rejoiced in God my Saviour'' (Luke 1:47): ``Because the
spirit of the Virgin rejoices in the eternal Godhead of the same Jesus, (i. e.
the Saviour,) whose flesh is formed in the womb by a temporal conception.''
The Collect of Ember Friday sums up the whole day's urgency in one line:
\textit{Excita, quaesumus, Domine, potentiam tuam, et veni}, asking that those
who trust in his kindness and love may be delivered speedily from every
adversity.

\subsection{Ember Saturday: The Vigil at St.~Peter's}

On Saturday the Roman church gathers at St.~Peter's, where the bishop of Rome
traditionally conferred sacred orders after the overnight fast. The Introit
cries out with Psalm 79~[80]:4, \textit{Veni, et ostende nobis}, ``Convert us,
O God: and shew us thy face, and we shall be saved,'' and the opening Collect
faces human weakness plainly:

\begin{massbox}{Opening Collect of Ember Saturday}
\textbf{Latin.} Deus, qui conspicis, quia ex nostra pravitate affligimur:
concede propitius; ut ex tua visitatione consolemur: Qui vivis et regnas cum
Deo Patre in unitate.

\textbf{English.} O God, Who dost behold how we are afflicted by reason of our
wickedness, mercifully grant that we may be consoled by reason of Thy
visitation. Who livest.
\end{massbox}

Because Ember Saturday preserves the architecture of an ancient nocturnal vigil,
the 1962 Missal prints two complete forms of its liturgy. In the longer form
(\textit{forma missae longior}), five Old Testament lessons unfold before the
Epistle, each followed by a chant and a Collect: four passages from Isaiah
(19:20--22; 35:1--7, where the desert blooms like the lily and the blind see;
40:9--11, the shepherd gathering lambs into his bosom; and 45:1--8, ending once
more with \textit{Rorate caeli}) culminate in Daniel 3:47--51 and the hymn of
the Three Young Men walking unharmed in the fiery furnace, followed by the
prayer \textit{Deus, qui tribus pueris}, asking that ``the flames of our sins
may not consume us Thy servants.'' In the shorter form, after the opening
Collect (\textit{Deus, qui conspicis}), one lesson from Isaiah (19:20--22) and
one additional Collect (\textit{Concede, quaesumus}) lead directly into the
Epistle.

Both forms then read 2 Thessalonians 2:1--8---warning Christians not to be
shaken by rumors about the Day of the Lord---and pass through the Tract
\textit{Qui regis Israel} (``Give ear, O thou that rulest Israel'') to the
Gospel of Luke 3:1--6, where John the Baptist steps into the fifteenth year of
Tiberius Caesar crying, ``Prepare ye the way of the Lord, make straight his
paths.'' At Matins on Ember Saturday, St.~Gregory the Great's Twentieth Homily
on the Gospels opens that very passage, observing why Luke names one Roman
emperor over the wide world and four petty rulers over small Judea: ``in the
whole heathen Commonwealth we find the title of one Emperor, but in the small
kingdom of Jud\ae a are mentioned four masters,'' foreshadowing how the nations
would be gathered into one while divided Judea stood in need of its Redeemer.

Whenever 21 December---the second-class feast of St.~Thomas the Apostle---falls
on Ember Friday (as in 2029 and 2035) or Ember Saturday (as in 2024 and 2030),
the universal feast of the Apostle takes precedence at Mass and in the Office,
while the Ember Day receives a privileged commemoration under the 1960 rubrics.

Even though modern canon law no longer binds Latin Catholics to fast on the
Winter Ember Days (Section~\ref{sec:penance-law}), these three days remain the
natural heart of Advent's inward turning. Keeping a simple voluntary fast on
Wednesday, Friday, or Saturday---putting what is saved aside for a neighbor in
need, reciting the \textit{Angelus} slowly on Wednesday and the
\textit{Magnificat} on Friday (Appendix~\ref{app:prayers}), and praying for
priests and seminarians on Saturday---joins a household directly to St.~Leo's
ancient harvest of mercy.
